\section{Further simulation studies}\label{supp:sim-more}

\subsection{Data-generating process}\label{sec:sim-dgp}

The simulator mimics the structure of the ACLED grid. The domain is a
$40 \times 30$ block of $0.5^\circ$ cells between $10^\circ$E and $30^\circ$E and
between $5^\circ$S and $10^\circ$N (1{,}200 cells of about 55\,km), with 104 training
weeks, 52 test weeks and 12 burn-in weeks. Countries are the Voronoi cells of randomly
placed seeds (14 countries). Four biomes are latitude bands with boundaries
perturbed by a smooth random field, ordered from ``cold'' (biome 1) to ``tropical''
(biome 4). Population is a standardised Mat\'ern field (range 250\,km) plus white
noise. Temperature and precipitation anomalies are standardised AR(1) processes in time
(coefficients 0.6 and 0.3) of Mat\'ern fields with ranges of 800 and 500\,km. A raw
temperature is also generated, as a climatology $30 - 1.2(\varphi - \varphi_{\min})$
in latitude $\varphi$, which rises from biome 1 to biome 4, plus a seasonal cycle
$(1 + 0.15|\varphi|)\cos(2\pi t/52)$ and 2.5 times the anomaly. The climatology and
the seasonal cycle have no effect in the truth. Three national covariates
$Z_0, Z_1, Z_2$ are standard normal across countries. $Z_1$ and $Z_2$ have no effect.
Outcomes are generated sequentially from
\begin{equation}\label{eq:sim}
\logit\pi_{st} = \beta_0 + \beta_{\mathrm{pop}}\mathrm{pop}_s
+ \delta\log\Bigl(1 + \textstyle\sum_{l=1}^{4} y_{s,t-l}\Bigr)
+ \beta(s)\,\mathrm{temp}_{st} + \beta_{\mathrm{prec}}\mathrm{prec}_{st}
+ \alpha Z_{0,c(s)} + u_{c(s)} + w(s) + v(s,t),
\end{equation}
with $\beta(s) = \gamma_{b(s)} + \eta(s)$ and $u_c \sim \Normal(0, 0.3^2)$. To avoid
generating data from the model used for fitting, $w$, $\eta$ and the innovations of
$v$ are drawn from Gaussian processes with the continuous Mat\'ern covariance
($\nu = 1$) in great-circle distance, evaluated at the cell centres. They are not drawn
from the lattice GMRF of Section~\ref{sec:methods-spde}. Unless stated otherwise,
$\beta_0 = -4$, $\beta_{\mathrm{pop}} = 0.5$, $\delta = 0.6$,
$\bgamma = (-0.3, -0.1, 0.1, 0.3)$, $\eta \equiv 0$, $\beta_{\mathrm{prec}} = -0.15$,
$\alpha = 0$, $\sigma_w = 1$, $\rho_w = 400$\,km and $v \equiv 0$. The resulting
prevalence of event weeks is about 4\%.

The default parameters are not calibrated to the application, and all covariate
effects in \eqref{eq:sim} are linear on the logit scale.

With $R$ replicates, a proportion $p$ (coverage, rejection rate) has Monte Carlo
standard error $\{p(1-p)/R\}^{1/2}$, at most 0.09 for $R = 30$ and 0.11 for $R = 20$, and
a mean estimate has Monte Carlo standard error equal to its empirical SD divided by
$R^{1/2}$. We give these standard errors where a comparison depends on them.

\subsection{S0: accuracy of the field representation and of the Laplace engine}\label{sec:sim-engine}

We draw $w$ on the $40 \times 30$ domain from the continuous Mat\'ern process with
$\rho_w = 400$\,km and $\sigma_w = 1$, and Bernoulli outcomes from
$\logit\pi_{st} = \beta_0 + 0.5\,x_{st} + w(s)$ with $x_{st} \sim \Normal(0,1)$, in two
settings with 50 replicates each. In the \emph{common} setting, $\beta_0 = -2$ with 40
weeks (prevalence about 16\%). In the \emph{rare} setting, $\beta_0 = -4.5$ with 104
weeks (prevalence about 2\%, as in the data). We also report the marginal variance of
the lattice field itself, from 2{,}000 exact draws, in interior cells and at the edge
of the domain. Table~\ref{tab:s0} summarises the results.

The lattice field has marginal variance 1.07 in cells at least eight cells (about
440\,km) from the edge of the domain, against a nominal 1, and 2.02 at the edge. This is
the doubling induced by the reflecting boundary that motivates the barrier model of
Remark~\ref{rem:boundary}. Although the data are not generated from the lattice model,
the range is recovered with little bias in both settings ($\hat\rho_w = 420$ and
409\,km against 400). The standard deviation is underestimated by about 12\%
($\hat\sigma_w = 0.88$ against 1). The bias is the same at prevalences of 16\% and
2\%, so it is not mainly due to the limited information in rare binary outcomes. A
plausible cause is the mismatch between the lattice model and the continuous process,
including the inflated variance near the boundary, since only about 28\% of the cells
are at least eight cells from the edge. We have not tested this explanation. The
posterior mean of $w$ has correlation 0.95 with the truth in the common setting and
0.91 in the rare one. Plug-in 95\% intervals cover $b_1$ in 88\% of replicates in the
common setting and 96\% in the rare one (Monte Carlo SE about 4.6\% and 2.8\%). They
cover $\beta_0 + w(s)$ for 94.8\% and 94.3\% of cells. Ignoring the uncertainty in
$\btheta$ therefore costs little here, including at the prevalence of the data. This
holds for the coefficient of a covariate that is independent in space and time and
independent of $w$, which is the easiest case. S0 does not check spatially smooth
covariates that are correlated with the field, and S1 shows that the plug-in standard
error of $\delta$ is too small when a space-time field is fitted. A fit takes 2--4
seconds on one core.

\begin{table}
\caption{S0: accuracy of the finite-volume SPDE field and of the empirical Bayes
Laplace approximation, with fields drawn from the continuous Mat\'ern process (50
replicates per setting). Mean over replicates, with the SD across replicates in
parentheses, except for the coverage of $b_1$, where it is the Monte Carlo standard
error.}\label{tab:s0}
\centering\small
\begin{tabular}{@{}lcc@{}}
\toprule
Quantity & common ($b_0=-2$, 50 reps) & rare ($b_0=-4.5$, 50 reps) \\
\midrule
Prevalence of event weeks & 0.161 (0.030) & 0.019 (0.005) \\
$\hat\rho_w$ (km) [truth 400] & 420 (40) & 409 (51) \\
$\hat\sigma_w$ [truth 1] & 0.884 (0.064) & 0.875 (0.075) \\
$\hat b_1$ [truth 0.5] & 0.494 (0.016) & 0.502 (0.021) \\
Coverage of 95\% interval for $b_1$ & 0.88 (0.05) & 0.96 (0.03) \\
Coverage of 95\% intervals for $b_0 + w(s)$ (share of cells) & 0.948 (0.009) & 0.943 (0.017) \\
$\mathrm{cor}(\hat w, w)$ & 0.954 (0.008) & 0.914 (0.018) \\
Fit time (s) & 2.4 (0.7) & 3.4 (0.8) \\
Marginal variance of $w$, interior cells [nominal 1] & 1.07 & -- \\
Marginal variance of $w$, edge cells & 2.02 & -- \\
\bottomrule
\end{tabular}

\end{table}

\subsection{S1: state dependence versus persistent heterogeneity}\label{sec:sim-persistence}

Three data-generating processes on a smaller $16 \times 14$ domain (224 cells, 6
countries, 104 training weeks, 20 replicates each) produce autocorrelation in $Y_{st}$
for different reasons:
\begin{itemize}
\item \emph{state dependence}: $\delta = 1$ and a weak field ($\sigma_w = 0.3$);
\item \emph{heterogeneity}: no state dependence ($\delta = 0$), a strong static field
($\sigma_w = 1.3$) and a space-time AR(1) field ($\sigma_v = 0.7$, $\phi = 0.9$ per
week, $\rho_v = 300$\,km);
\item \emph{mixed}: $\delta = 0.5$, $\sigma_w = 1$ and a space-time field with
$\sigma_v = 0.5$ and $\phi = 0.9$.
\end{itemize}
Two further truths repeat the heterogeneity and mixed scenarios with a more persistent
space-time field ($\phi = 0.97$ per week). We fit model~\eqref{eq:model} with lags only,
with lags and $w$, and with lags, $w$ and $v$. The space-time field $v$ is resolved at
weekly steps ($\Delta = 1$, the resolution of the data-generating process), at four-week
steps ($\Delta = 4$) and at eight-week steps ($\Delta = 8$, the resolution of the
regional state-dependence fits in the application). Table~\ref{tab:s3} shows the three
main truths at $\Delta = 1$ and $\Delta = 4$. The $\Delta = 8$ fits, the two
persistent truths and the fit times are reported in the text only.

Under state dependence, ignoring the latent field inflates $\hat\delta$ by about a
quarter (1.26; coverage 0.60), and all latent specifications are approximately unbiased
(0.98--1.03) with coverage of 0.90--0.95. Under heterogeneity the contrast is much
sharper. With lags only, $\hat\delta = 1.61$ although the true value is zero. Adding the
static field reduces the estimate to 0.46, but no interval covers zero: slowly varying
episodes act as state dependence. Only the models with the space-time field bring the
estimate close to zero: $\hat\delta$ is 0.05 with four-week steps and 0.01 with weekly
steps, both within Monte Carlo error of zero (Monte Carlo SE about 0.035). The mixed
scenario behaves in the same way. The static field leaves an upward bias (0.78 against
0.5), which the space-time field removes (0.55 and 0.54, Monte Carlo SE about 0.04).

The standard errors of $\hat\delta$ are too small when $v$ is fitted. In the
heterogeneity and mixed scenarios the empirical SD exceeds the mean standard error by
27--43\% (ratios 1.27--1.43), and coverage is 0.80--0.90 (Monte Carlo SE 0.07--0.09).
In the state-dependence scenario, where the truth has no $v$, the ratio is
0.92--1.02 for the four latent models. With 20
replicates each ratio has a Monte Carlo SE of about 0.2, so the size of the shortfall is
imprecise, but it appears in all twelve cells with $v$ in the truth and in the fitted
model (ratios 1.19--1.66, including the $\phi = 0.97$ truths and $\Delta = 8$, with
coverage 0.80--0.90). Plug-in intervals for $\delta$ from the space-time model are
therefore too short, by a factor of roughly 1.2 to 1.7 in these designs. The ratio is
not stable enough to justify a single correction factor. In the application, plug-in
standard errors from models with a space-time field should be read as too small, and
an estimate about two plug-in standard errors from zero should not be read as clearly
different from zero.

Resolving $v$ at four-week steps instead of weekly steps changes $\hat\delta$ little in
these scenarios, and it is four to six times faster (ratios of mean fit time 3.9 to
5.6). Paired by replicate, the four-week
estimate is larger by 0.034 (SE 0.006) under heterogeneity and by 0.010 (SE 0.008) in
the mixed scenario. The space-time field has an AR(1) coefficient of 0.9 per week, so it
changes slowly relative to the four-week history window. This argument is weaker at
eight-week steps, since $0.9^8 \approx 0.43$. With
$\Delta = 8$ the estimate is still unbiased under state dependence (0.98, coverage
0.95), but under heterogeneity it is 0.10 (empirical SD 0.14, mean SE 0.12, coverage
0.80), about three Monte Carlo SE above zero. Paired by replicate, it exceeds the weekly
estimate by 0.086 (SE 0.012). In the mixed scenario it is 0.57 against 0.5 (coverage
0.85; 0.031 above the weekly estimate, SE 0.011). Coarse steps therefore leave part of
the episodes to be absorbed by the lag. The bias is smaller when the episodes are more
persistent. With $\phi = 0.97$, $\hat\delta$ at $\Delta = 8$ is 0.04 under
heterogeneity (coverage 0.90) and 0.50 in the mixed scenario (coverage 0.85), and it
exceeds the weekly estimate by 0.050 (SE 0.007) and 0.024 (SE 0.006). In these slow
truths the static field alone leaves less bias than before (0.38 and 0.71, against 0.46
and 0.78), and the four-week and weekly estimates are within 0.03 of the truth. A fit
at $\Delta = 8$ takes 23--27 seconds, against 84--101 at $\Delta = 4$ and 350--474
at $\Delta = 1$. Faster space-time dynamics ($\phi < 0.9$) are
not covered, and neither are the 13-week steps of the joint and space-time fits in
the application.

\begin{table}
\caption{S1: estimation of the state-dependence parameter $\delta$ (20 replicates per
scenario).}\label{tab:s3}
\centering\small
\begin{tabular}{@{}lcccc@{}}
\toprule
Estimator & Mean $\hat\delta$ & Emp.\ SD & Mean SE & Coverage \\
\midrule
\multicolumn{5}{@{}l}{\emph{State dependence ($\delta = 1$)}} \\
\quad Lags only & 1.263 & 0.107 & 0.136 & 0.60 \\
\quad Lags + $w$ & 1.032 & 0.131 & 0.139 & 0.90 \\
\quad Lags + $w$ + $v$ (4-week steps) & 0.990 & 0.135 & 0.140 & 0.95 \\
\quad Lags + $w$ + $v$ (weekly steps) & 0.997 & 0.143 & 0.140 & 0.95 \\
\quad XGBoost: share of mean $|$SHAP$|$ for lags & 0.108 & 0.034 & -- & -- \\
\multicolumn{5}{@{}l}{\emph{Heterogeneity ($\delta = 0$)}} \\
\quad Lags only & 1.614 & 0.386 & 0.099 & 0.00 \\
\quad Lags + $w$ & 0.464 & 0.093 & 0.111 & 0.00 \\
\quad Lags + $w$ + $v$ (4-week steps) & 0.045 & 0.152 & 0.119 & 0.90 \\
\quad Lags + $w$ + $v$ (weekly steps) & 0.011 & 0.162 & 0.120 & 0.90 \\
\quad XGBoost: share of mean $|$SHAP$|$ for lags & 0.094 & 0.038 & -- & -- \\
\multicolumn{5}{@{}l}{\emph{Mixed ($\delta = 0.5$)}} \\
\quad Lags only & 1.659 & 0.351 & 0.106 & 0.00 \\
\quad Lags + $w$ & 0.782 & 0.111 & 0.115 & 0.25 \\
\quad Lags + $w$ + $v$ (4-week steps) & 0.545 & 0.152 & 0.120 & 0.85 \\
\quad Lags + $w$ + $v$ (weekly steps) & 0.535 & 0.172 & 0.120 & 0.80 \\
\quad XGBoost: share of mean $|$SHAP$|$ for lags & 0.116 & 0.044 & -- & -- \\
\bottomrule
\end{tabular}

\end{table}
